% PDF-to-TeX transcription by the assistant; not the original downloaded file.
\begin{theorem}[Lecture 30, Theorem 1 (Strichartz)]
If $u$ solves $\partial_tu=i\Delta u$ and $u(x,0)=u_0(x)$,
\[\|e^{it\Delta}u_0\|_{L^\sigma_{x,t}}\lesssim\|u_0\|_{L^2_x},\tag{2}\]
where $\sigma=2(d+2)/d$.
\end{theorem}
\paragraph{Preparatory estimates (Lecture 29)}
First, the $L^2$ norm of a solution is preserved in time:
\[\|e^{it\Delta}u_0\|_{L^2_x}=\|u_0\|_{L^2_x}.\]
Second, solutions of the Schrodinger equation obey an $L^\infty$ decay estimate:
\[\|e^{it\Delta}u_0\|_{L^\infty_x}\lesssim|t|^{-d/2}\|u_0\|_{L^1_x}.\]
We have two estimates---the conservation of $L^2$ and the decay estimate. Now that we have proven the interpolation theorem, we can interpolate between these two estimates.
\begin{proposition}[Lecture 29, Proposition 2]
For any $0\leq\theta\leq1$, define $p$ by $1/p=(1-\theta)\cdot1/2$, and let $p'$ be the dual exponent. Then we have the following inequality.
\[\|e^{it\Delta}u_0\|_{L^p_x}\lesssim t^{-\frac d2\cdot\theta}\|u_0\|_{L^{p'}_x}.\]
\end{proposition}
\paragraph{The inhomogeneous Schrodinger equation (Lecture 29, \S2)}
We consider the inhomogeneous Schrodinger equation $\partial_tu=i\Delta u+F$. Here $u$ and $F$ are both functions of $x$ and $t$. We will write $F_t(x)$ for $F(x,t)$. Similarly, we will write $u_t(x)$ for $u(x,t)$.
\begin{proposition}[Lecture 29, Proposition 3 (Duhamel formula)]
If $F\in C^\infty_{\mathrm{comp}}(\R^d\times\R)$, then the following function $u$ solves the inhomogeneous Schrodinger equation:
\[u_t=\int_{-\infty}^te^{i(t-s)\Delta}F_s\,ds.\]
Moreover, the function $u(x,t)$ vanishes at all times $t$ ``before'' the support of $F$.
\end{proposition}
\begin{proof}
The last claim is easy to check. Suppose that $F$ is supported on $\R^d\times[T_1,T_2]$. If $t<T_1$, then $F_s=0$ for all $s\in[-\infty,t]$, and so $u_t=0$.

Recall that $e^{it\Delta}u_0$ solves the Schrodinger equation:
\[\partial_t(e^{it\Delta}u_0)=i\Delta(e^{it\Delta}u_0).\]
So taking the time derivative of $u_t$, we get
\begin{align*}
\partial_tu_t&=e^{i(t-s)\Delta}F_s\big|_{s=t}+\int_{-\infty}^t\partial_t(e^{i(t-s)\Delta}F_s)\,ds\\
&=F_t+\int_{-\infty}^ti\Delta(e^{i(t-s)\Delta}F_s)\,ds=F_t+i\Delta u_t.
\end{align*}
\end{proof}
\begin{theorem}[Lecture 29, Theorem 4 (Also Strichartz)]
Suppose that $u$ obeys the inhomogeneous Strichartz equation $\partial_tu=i\Delta u+F$, and that $u$ vanishes at times before the support of $F$. Let $s$ be the Strichartz exponent $s=2(d+2)/d$ as above, and let $s'$ be its dual exponent. Then
\[\|u\|_{L^s_{x,t}}\lesssim\|F\|_{L^{s'}_{x,t}}.\]
\end{theorem}
\begin{proof}
We will use the Duhamel formula, and the $L^p$ estimates in Proposition 2. For any $p$, we have
\[\|u\|_{L^p_{x,t}}^p=\int_\R\|u_t\|_{L^p_x}^p\,dx.\]
By Duhamel's formula and Minkowski's inequality,
\[\|u_t\|_{L^p_x}=\left\|\int_{-\infty}^te^{i(t-s)\Delta}F_s\,ds\right\|_{L^p_x}\leq\int_{-\infty}^t\|e^{i(t-s)\Delta}F_s\|_{L^p_x}\,ds.\]
As in Proposition 2, let's suppose that $1/p=(1-\theta)\cdot1/2$. Applying Proposition 2, we see that
\[\|u_t\|_{L^p_x}\lesssim\int_{-\infty}^t(t-s)^{-\frac d2\theta}\|F_s\|_{L^{p'}_x}\,ds.\tag{1}\]
The right-hand side is a convolution which is a little hard to see with all the notation. We let $g(s)=\|F_s\|_{L^{p'}_x}$, we let $h(s)=\|u_s\|_{L^p_x}$, and we let $\alpha=\frac d2\theta$. Then the last equation gives
\[h(t)\leq g*|t|^{-\alpha}(t).\tag{2}\]
Note that $\|h\|_p=\|u\|_{L^p_{x,t}}$ and $\|g\|_{p'}=\|F\|_{L^{p'}_{x,t}}$. By Hardy-Littlewood-Sobolev, and equation (2), we know that $\|h\|_r\lesssim\|g\|_q$ provided that
\[\frac1r+1=\frac1q+\alpha.\]
In particular, $\|h\|_p\lesssim\|g\|_{p'}$ as long as
\[\frac1p+1=\frac{p-1}{p}+\frac d2\cdot\theta.\]
When we plug in $p=s$ and find the corresponding $\theta$, this equation is satisfied, and so we get $\|u\|_{L^s_{x,t}}\lesssim\|F\|_{L^{s'}_{x,t}}$ as desired. We do the computation with $s$ and $\theta$ here in the notes for completeness, although I'm not sure if it's illuminating enough to include in the lecture.

Recall that $p$ and $\theta$ are related by $1/p=(1-\theta)/2$, which yields $\theta=1-2/p=(p-2)/p$. Plugging for $\theta$ in the last equation, we get
\[\frac1p+1=\frac{p-1}{p}+\frac{d(p-2)}{2p}.\]
Multiplying through by $2p$, we get
\[2+2p=2(p-1)+d(p-2),\qquad4=d(p-2),\qquad p=\frac4d+2=\frac{2(d+2)}d=s.\]
\end{proof}
\paragraph{Unitarity and the homogeneous estimate (Lecture 30)}
\begin{lemma}[Lecture 30, Lemma 2]
$\langle e^{it\Delta}f,g\rangle_{\R^d}=\langle f,e^{-it\Delta}g\rangle_{\R^d}$.
\end{lemma}
\begin{proof}
By the definition of the $L^2$-inner product on $\R^d$ and Plancherel's theorem:
\begin{align*}
\langle e^{it\Delta}f,g\rangle_{\R^d}&=\int_{\R^d}e^{it\Delta}f\,\overline g
=\int_{\R^d}e^{it(2\pi i\omega)^2}\widehat f\,\overline{\widehat g}\\
&=\int_{\R^d}\widehat f\,\overline{e^{-it(2\pi i\omega)^2}\widehat g}
=\int_{\R^d}f\,\overline{e^{-it\Delta}g}=\langle f,e^{-it\Delta}g\rangle_{\R^d}.
\end{align*}
\end{proof}
With this unitarity we may now proceed with the proof of Strichartz:
\begin{proof}[Proof of Lecture 30, Theorem 1]
By duality,
\[\|e^{it\Delta}u_0\|_{L^\sigma_{x,t}}=\sup_{\|F\|_{L^{\sigma'}_{x,t}}=1}\int_{\R^d\times\R}e^{it\Delta}u_0\overline F,\]
where $\sigma'$ is the dual exponent of $\sigma$ satisfying $1/\sigma+1/\sigma'=1$. By Lemma 2,
\begin{align*}
\sup_{\|F\|_{L^{\sigma'}_{x,t}}=1}\int_{\R^d\times\R}e^{it\Delta}u_0\overline F
&=\sup\int_\R\langle e^{it\Delta}u_0,F_t\rangle\,dt\\
&=\sup\int_\R\langle u_0,e^{-it\Delta}F_t\rangle\,dt
=\sup\left\langle u_0,\int e^{-it\Delta}F_t\,dt\right\rangle.
\end{align*}
Hence by Cauchy-Schwarz:
\[\|e^{it\Delta}u_0\|_{L^\sigma_{x,t}}\leq\|u_0\|_{L^2_x}\sup_F\left\|\int e^{-it\Delta}F_t\,dt\right\|_{L^2_x}.\]
It therefore suffices to check that
\[\sup_{\|F\|_{L^{\sigma'}_{x,t}}=1}\left\|\int e^{-it\Delta}F_t\,dt\right\|_{L^2_x}\lesssim1.\]
We prove this separately as its own lemma:
\end{proof}
\begin{lemma}[Lecture 30, Lemma 3]
$\left\|\int e^{-it\Delta}F_t\,dt\right\|_{L^2_x}\lesssim\|F\|_{L^{\sigma'}_{x,t}}$.
\end{lemma}
\begin{proof}
\begin{align*}
\left\|\int e^{-it\Delta}F_t\,dt\right\|_{L^2_x}^2
&=\left\langle\int_\R e^{-it\Delta}F_t\,dt,\int_\R e^{-is\Delta}F_s\,ds\right\rangle\\
&=\iint_{\R^2}\langle e^{-it\Delta}F_t,e^{-is\Delta}F_s\rangle\,dt\,ds\\
&=\iint_{\R^2}\langle F_t,e^{i(t-s)\Delta}F_s\rangle\,dt\,ds.
\end{align*}
This expression effectively measures the interaction between all pairs $(t,s)$, as highlighted earlier. Now if
\[G(x,t)=\int_\R e^{i(t-s)\Delta}F_s(x)\,ds,\]
we may write
\begin{align*}
\left\|\int e^{-it\Delta}F_t\,dt\right\|_{L^2_x}^2
&=\iint_{\R^2}\langle F_t,e^{i(t-s)\Delta}F_s\rangle\,dt\,ds\\
&=\int_\R\left\langle F_t,\int_\R e^{i(t-s)\Delta}F_s\,ds\right\rangle\,dt
=\int_{\R^d\times\R}F\overline G.
\end{align*}
By H\"older,
\[\left\|\int e^{-it\Delta}F_t\,dt\right\|_{L^2_x}^2\leq\|F\|_{L^{\sigma'}_{x,t}}\|\overline G\|_{L^\sigma_{x,t}}.\]
Finally, by the Duhamel bound derived in the previous lecture, $\|\overline G\|_{L^\sigma_{x,t}}\lesssim\|F\|_{L^{\sigma'}_{x,t}}$. Hence
\[\left\|\int e^{-it\Delta}F_t\,dt\right\|_{L^2_x}^2\lesssim\|F\|_{L^{\sigma'}_{x,t}}^2.\]
\end{proof}
